\documentclass{article}
\input{../../../lib.sty}

\title{\texttt{fn expected\_nb\_mean\_from\_x\_hi}}
\author{Michael Shoemate}

\begin{document}
\maketitle

Proves soundness of \texttt{fn expected\_nb\_mean\_from\_x\_hi} in \texttt{mod.rs}.

\section{Hoare Triple}
\subsection*{Precondition}
\texttt{eta} $\ge 0$ and \texttt{x} $> 0$.

\subsection*{Pseudocode}
\lstinputlisting[language=Python,firstline=2,escapechar=|]{./pseudocode/expected_nb_mean_from_x_hi.py}

\subsection*{Postcondition}
\begin{theorem}
\label{postcondition}
For any inputs satisfying the precondition,
\texttt{expected\_nb\_mean\_from\_x\_hi} either returns \texttt{Err(e)},
or returns \texttt{Ok(out)} where \texttt{out} is an upper enclosure of the Definition 1
expected repetition count.
\end{theorem}

\section{Proof}
Assume the precondition, let $q = e^{-x}$ and $\gamma = 1 - e^{-x}$.
Since $x > 0$, we have $q \in (0,1)$ and $\gamma \in (0,1)$.

\begin{lemma}
\label{err-e}
\texttt{expected\_nb\_mean\_from\_x\_hi} only returns \texttt{Err(e)} if one of the fallible calls
on lines \ref{line:q_hi}, \ref{line:gamma_lo}, \ref{line:interval}, \ref{line:gamma_eta_hi},
or the unlabeled fallible arithmetic calls in the taken branch returns \texttt{Err(e)}.
\end{lemma}

\begin{proof}
Every branch return is either an explicit \texttt{float("inf")} or the result of propagating one
of the listed fallible helper calls.
\end{proof}

\begin{lemma}
\label{log-branch}
If \texttt{eta = 0} and \texttt{expected\_nb\_mean\_from\_x\_hi} returns \texttt{Ok(out)}, then
\[
\texttt{out} \ge \frac{e^{-x}}{\gamma \log(1/\gamma)}.
\]
\end{lemma}

\begin{proof}
By the proof definitions of \texttt{inf\_exp} and \texttt{gamma\_lo\_from\_x},
\[
\texttt{q\_hi} \ge q, \qquad \texttt{gamma\_lo} \le \gamma.
\]
If line \ref{line:interval} succeeds, then \texttt{gamma\_lo} lies in $(0,1)$.
By the proof of \texttt{log\_inv\_gamma\_lo\_from\_x},
\[
\texttt{log\_inv\_gamma\_lo} \le \log(1/\gamma).
\]
Therefore the rounded multiplication on line \ref{line:logarithmic_branch} computes a lower bound
on the true denominator, so
\[
\texttt{denom\_lo} \le \gamma \log(1/\gamma).
\]
If \texttt{denom\_lo $\le 0$}, the function returns $+\infty$, which is an upper bound.
Otherwise \texttt{denom\_lo} is a positive lower bound on the true denominator, so upward-rounded
division yields
\[
\texttt{out} \ge \frac{q}{\gamma \log(1/\gamma)} = \frac{e^{-x}}{\gamma \log(1/\gamma)}.
\]
\end{proof}

\begin{lemma}
\label{nb-branch}
If \texttt{eta > 0} and \texttt{expected\_nb\_mean\_from\_x\_hi} returns \texttt{Ok(out)}, then
\[
\texttt{out} \ge \frac{\texttt{eta} \cdot e^{-x}}{\gamma(1 - \gamma^{\texttt{eta}})}.
\]
\end{lemma}

\begin{proof}
As above, \texttt{q\_hi} $\ge q$ and \texttt{gamma\_lo} $\le \gamma$, with \texttt{gamma\_lo} in
$(0,1)$ after line \ref{line:interval}.
By the proof of \texttt{gamma\_pow\_eta\_hi\_from\_x},
\[
\texttt{gamma\_eta\_hi} \ge \gamma^{\texttt{eta}}.
\]
Therefore
\[
\texttt{one\_minus\_gamma\_eta\_lo} \le 1 - \gamma^{\texttt{eta}}.
\]
If \texttt{one\_minus\_gamma\_eta\_lo $\le 0$} or \texttt{denom\_lo $\le 0$}, the function returns
$+\infty$, which is an upper bound.
Otherwise both lower bounds are positive, so
\[
\texttt{denom\_lo} \le \gamma(1 - \gamma^{\texttt{eta}})
\quad\text{and}\quad
\texttt{num\_hi} \ge \texttt{eta} \cdot e^{-x}.
\]
Upward-rounded division then gives the claimed upper bound.
\end{proof}

\begin{proof}[Proof of Theorem~\ref{postcondition}]
Holds by Lemma~\ref{err-e}, Lemma~\ref{log-branch}, and Lemma~\ref{nb-branch}.
\end{proof}

\end{document}
